---
title: "Introduction to Mean, Median and Mode"
description: "Mean, Median and Mode – Measures of Central Tendency"
pubDate: 2024-01-15  # Add this - use current date
author: "Your Name"   # Add your name
layout: ../../layouts/BlogPost.astro  # This might be needed
draft: false          # Explicitly set to false
published: true       # Explicitly set to true
math: true
tags:
  - statistics
  - mathematics
  - data-analysis
---

## **Introduction**

All three are measures of central tendency.

**Mean** is just the average of the data.

**Median** is the middle value of the data when it is arranged in ascending order.  
If there is an even number of data points, the median is the average of the two middle values.  
For example, if the data is:

$$
2,\; 3,\; 4,\; 2,\; 3,\; 4,\; 4
$$

then the median is:

$$
\frac{3 + 4}{2} = 3.5
$$

**Mode** is the number that appears most frequently in a data set.  
In the above example, the mode is:

$$
4
$$

It is important to note that a data set may have more than one mode if multiple values occur with the same highest frequency.

the end 